\section{Falsification and Simplification Criteria}

ARC is intended to become smaller when evidence does not support added structure.

\subsection{Core temporal sequence}

The current H1 pilot separates two coefficient-level claims from structural and predictive evidence. H1-P1 tests the desired-state coefficient in M2 after adjustment for perceived state, context, and action history. H1-P2 tests the signed-discrepancy coefficient in M3. Both use the locked one-sided SESOI tests with Holm family-wise control. H1-P2 alone does not establish a contribution beyond perceived state: M3 does not include a separate perceived-state coefficient, and the difference score contains perceived-state information. Even joint support does not uniquely identify discrepancy regulation as a psychological mechanism.

The target component should be simplified if adequately powered studies repeatedly find negligible within-person target variation and no incremental target information. A nonsignificant coefficient or an imprecise prediction difference in one pilot is insufficient evidence of absence. If target drift is substantial, it must be modeled explicitly rather than treated as error. More consistent reverse or simultaneous associations would challenge the proposed temporal ordering.

The difference-score representation is a testable restriction. For M2, report \(d_{\mathrm{sym}}=\beta_{r^\ast}+\beta_{\hat r}\) with the prespecified 90\% confidence interval and tolerance \(\delta_{\mathrm{sym}}=1.0\) on the coefficient-per-10-point scale. An interval wholly inside \([-1,1]\) supports practical compatibility with the restriction; an interval wholly outside it on one side provides evidence against it; overlap is inconclusive. Compatibility within a tolerance does not establish exact equality or a unique mechanism.

For M4, report the positive and negative discrepancy slopes and \(d_{\mathrm{asym}}=\beta_{+}+\beta_{-}\), where \(\beta_{-}\) multiplies \(\max(-e,0)\). Report the estimate and 95\% confidence interval descriptively. The focused validation did not justify a confirmatory asymmetry tolerance at \(N=84\). Therefore, this pilot does not classify M4 as practically symmetric or asymmetric, and an interval containing zero is not evidence of symmetry.

Participant-paired temporal-holdout RMSE differences and their uncertainty are secondary predictive evidence. They neither gate the co-primary tests nor impose a mandatory winning model. Reproducible evidence against the M2 coefficient restriction motivates separate state and target pathways; replicated asymmetric patterns with adequate precision motivate a prospectively calibrated asymmetric extension. Structural incompatibility with little predictive loss means that the simple representation may remain a useful approximation, not that its psychological mechanism has been established. Discordant or imprecise evidence must remain visible rather than be converted into a favorable model choice.

\subsection{Attachment moderators}

The dynamic role assigned to \Anx\ and \Avo\ should be removed if they do not reproducibly moderate discrepancy-to-action or related temporal parameters. Their established value as descriptive attachment dimensions would remain; only the ARC-specific dynamic interpretation would fail.

H2 is specifically weakened if attachment predicts only between-person mean levels of approach, distancing, or internal regulation but does not alter within-person discrepancy-to-action mappings. If moderation appears only in a signed action summary but not in its component channels, the channel decomposition should take precedence because opposite high-intensity behaviors can cancel in a signed score.

A null interaction in an underpowered pilot is not sufficient for rejection. H2 should be evaluated only in a design simulated for cross-level moderation and random-slope heterogeneity. Conversely, a statistically detectable interaction that does not replicate or improve held-out prediction should not be retained as an ARC policy parameter.

\subsection{Positive security candidate}

\Sec\ should remain outside the required ARC core unless a positive secure-base measure improves a preregistered primary recovery outcome beyond \Anx\ and \Avo, adds genuinely out-of-sample predictive value, and the increment is independently replicated or supported by a second preregistered design. A zero-order association between \Sec\ and recovery is insufficient evidence for H3 or for promotion to core status. A reduction in discrepancy that is explained only by retreat of the desired state is also insufficient evidence of successful recovery.

If \Sec\ adds information beyond \Anx/\Avo\ but loses that information after proximal partner support or observed dyadic coping is included, ARC should report attenuation conditional on support rather than infer mediation or automatically promote \Sec\ as an independent person-level core construct. Attenuation alone does not identify mediation: shared measurement, confounding, and correlated predictors are alternative explanations. A mediation claim requires a separately specified temporal causal model and defensible identification assumptions. If proximal support predicts recovery while \Sec\ does not, the recovery account should be simplified toward episode-level coregulation.

If \Sec\ predicts relationship quality or secure-base behavior but not the preregistered primary recovery outcome, the positive construct may remain theoretically relevant, but it should remain outside the ARC core and its placement in the dynamic recovery pathway should be removed. Secondary outcomes cannot be substituted post hoc for a failed primary promotion endpoint.

\subsection{Mentalized affectivity and flexibility}

\MentAff\ and \Flex\ should be excluded, merged with adjacent constructs, or redefined if they fail discriminant validity or do not add incremental information to the specific pathway assigned to them. ARC does not depend on preserving these moderators.

\subsection{Dynamic disorganization extension}

Dynamic indicators should not be described as attachment disorganization if they fail to improve held-out prediction of an independent disorganization criterion beyond \Anx, \Avo, their interaction, and simpler static summaries. High variability by itself is insufficient evidence of disorganization.

\subsection{Mechanistic non-uniqueness}

Even successful prediction of the ARC temporal sequence would not uniquely identify a latent computational mechanism. If P/PI/PID, state-space, active-inference, reinforcement-learning, or other models remain observationally equivalent, the defensible conclusion is about temporal prediction and feedback organization rather than a unique psychological algorithm. This identifiability problem is addressed separately in Manuscript 2.

